\chapter{总结与展望}\label{chap:conclusion}

\section{本文工作总结}

本文围绕宝可梦PVP对战策略优化这一核心问题，从属性覆盖、速度优化和队伍协同三个维度进行了系统研究，取得了以下成果：

\begin{enumerate}
    \item \textbf{属性覆盖面优化}。提出了融合双属性组合与特性修正的GLSH属性覆盖优化算法。实验表明，该方法使队伍的属性覆盖度达到91.3\%，相比基于使用率的选队方法提升12.4个百分点。引入双属性组合效应因子有效避免了单属性评估的覆盖度虚高问题。

    \item \textbf{速度线阈值预测}。构建了基于贝叶斯优化的速度线阈值预测模型。该模型通过高斯过程代理和期望提升采集函数，实现了对关键速度门槛的高精度估计（MAE = 3.2点速度值），为训练家的努力值分配提供了定量依据。

    \item \textbf{多目标队伍协同优化}。设计了基于NSGA-II的多目标遗传算法框架，同时优化属性覆盖度、速度配合度和攻防平衡度。综合胜率达到68.0\%，显著优于现有方法。

    \item \textbf{实验验证}。构建了包含500万场对战记录的数据集，从覆盖度、速度配合度和胜率三个维度对所提方法进行了全面评估。消融实验验证了各算法组件的独立贡献。
\end{enumerate}

\section{主要创新点}

本文的主要创新点可归纳如下：

\begin{enumerate}
    \item 在属性覆盖优化中首次同时考虑双属性组合效应和特性修正项，使覆盖度评估更加精确；
    \item 将贝叶斯优化引入速度线分析，实现了从经验判断到定量预测的转变；
    \item 提出面向队伍整体协同的多目标优化框架，协调了属性、速度和攻防三个维度的优化需求。
\end{enumerate}

\section{未来研究方向}

尽管本文在宝可梦PVP对战策略优化方面取得了一定成果，但仍存在以下值得进一步探索的方向：

\subsection{动态环境适应}

当前方法假设对战环境（Meta）是静态的，即宝可梦的使用率分布保持不变。然而，实际对战环境会随着赛季更新和玩家偏好的变化而动态演进。未来的研究可以引入在线学习机制，使队伍配置能够自动适应环境变化。

\subsection{双打对战协同建模}

本文的方法主要针对6v6单打对战设计。在4v4双打对战中，宝可梦之间的协同配合（如天气队、空间队、顺风队等战术体系）更为复杂。未来的工作可以将双打协同关系显式建模为图结构，通过图神经网络学习宝可梦之间的配合模式。

\subsection{对手行为预测}

当前的队伍优化方法假设对手的策略是静态的，未考虑对手的实时决策。将博弈论中的纳什均衡和对手建模（Opponent Modeling）引入队伍优化，有望进一步提升实战表现。

\subsection{大语言模型辅助策略生成}

大语言模型（LLM）在策略推理和知识表示方面展现出强大能力。利用LLM对宝可梦对战的文本描述（如Smogon战术分析、VGC赛后报告）进行语义理解，辅助生成可解释的战术建议，是值得探索的前沿方向。

\section{结语}

宝可梦对战作为一项兼具策略深度和娱乐性的竞技活动，其背后的优化问题涉及组合优化、博弈论、机器学习等多个学科领域。本文的工作只是这一交叉研究方向的初步探索，期待更多研究者加入，共同推动宝可梦对战策略研究的发展，为训练家们提供更科学、更系统的队伍配置决策支持。
